\DO{This list should be filled out fully and submitted together
  \emph{with your submissions until the final one!}} Insert it as a
\verb|\subsection*{Project report checklist}| right by the end of
\secref{intro}. In submitting the final and revised version you simply
skip this checklist.

\begin{itemize}
\item Titlepage, Abstract, Table of contents
  \begin{todolist}
  \item The title is the correct and final title
  \item All author's names are there, usually in alphabetic order as
    to the family name
  \item The abstract is informative and focuses on the problem
    approached and what you did
  \item The abstract does not contain undefined abbreviations,
    acronyms, references, equations, etc., unless you have good
    reasons
  \item The table of contents is up to date and makes a structured and
    clear impression
   \itemdone Checked items are marked out using \verb!\itemdone! and
   all checked items are actually taken care of!
  \end{todolist}
  
\item Structure
  \begin{todolist}
  \item The report makes sensible use of sections and subsections:
    avoid subsubsections and appendices unless you have good reasons
  \item In the case there is an appendix, it is explained why it is
    there in its initial paragraph
  \item The first section (introduction) explains what will come in the
    report and in what order (usually in the final paragraph)
  \item Each section starts with a few sentences explaining what will
    come in the section
  \end{todolist}

\item Contents
  \begin{todolist}
  \item Each section is made up of paragraphs which gives a coherent
    flow to a reader and concerns the topic of the section
  \item Each paragraph is made up of sentences which gives a coherent
    flow to a reader and concerns the topic of the paragraph
  \item Other than common knowledge or trivialities, no claims are made
    without stating the reasons or inserting a proper citation which
    supports the claim
  \item Equations are used when needed and are referenced when useful
  \item The report has been consistently spell-checked and proofread by
    all members in the project group
  \item You have received feedback from your supervisor(s) and you have
    made a serious attempt to adress all comments received
  \item Volume: the normal range is about 12--24 pages; in case the
    report is shorter, then there are reasons for this, and in case the
    report is substantially longer, then there are reasons for this
  \end{todolist}
  
\item References, Figures, Tables
  \begin{todolist}
  \item References are used throughout to support claims and direct
    the reader to useful sources of information
  \item The actual items in the reference list by the end uses
    \verb!\doi{}!, in some cases \verb!\url{}!, but \emph{never
      both}. \emph{Never} use references which involves a
    Google-search or similar. Similarly \emph{never} use references
    from ScienceDirect or similar pages as these URLs are not original
    ones!
  \item The use of ISSN/ISBN numbers have been avoided
    altogether. Also, generally, only years are listed, so not
    individual months.
  \item For all cited arXiv preprints, these have been checked in case
    they are actually published. In case they were published, this
    version is rahter the one cited.
  \item No figure is taken from other sources without either a written
    permission (stated) or under an open license (also stated)
  \item Captions for figures and tables explain what information is
    communicated
  \item Captions are a full short paragraph and hence ends with a
    period
  \item Figure titles are avoided; units are clearly desiplayed and
    the font for the axis is reasonably large
  \item All figures included in the report are meaningful and are
    meant to be there since they communicate information
  \item All figures and tables are discussed in the main text and are
    then also referenced
  \end{todolist}

\item Use of AI
  \begin{todolist}
  \item All AI-generated figures are marked out/cited as such, are
    included because they are meaningful, and accurately convey the
    information they are meant to
  \item The report generally and transparently details all uses of
    AI-tools, e.g., as a minium by explaining how it was used and to
    what extend (see the acknowledgment for an example)
  \item Text or code proposed or modified by an AI has been checked
    for correctness and that the meaning is as intended
  \item All members of the project team take full responsibility for
    the whole text
  \end{todolist}
\end{itemize}